% Options for packages loaded elsewhere
\PassOptionsToPackage{unicode}{hyperref}
\PassOptionsToPackage{hyphens}{url}
\PassOptionsToPackage{dvipsnames,svgnames,x11names}{xcolor}
\documentclass[
  11pt,
]{article}
\usepackage{xcolor}
\usepackage[margin=2.4cm]{geometry}
\usepackage{amsmath,amssymb}
\setcounter{secnumdepth}{-\maxdimen} % remove section numbering
\usepackage{iftex}
\ifPDFTeX
  \usepackage[T1]{fontenc}
  \usepackage[utf8]{inputenc}
  \usepackage{textcomp} % provide euro and other symbols
\else % if luatex or xetex
  \usepackage{unicode-math} % this also loads fontspec
  \defaultfontfeatures{Scale=MatchLowercase}
  \defaultfontfeatures[\rmfamily]{Ligatures=TeX,Scale=1}
\fi
\usepackage{lmodern}
\ifPDFTeX\else
  % xetex/luatex font selection
\fi
% Use upquote if available, for straight quotes in verbatim environments
\IfFileExists{upquote.sty}{\usepackage{upquote}}{}
\IfFileExists{microtype.sty}{% use microtype if available
  \usepackage[]{microtype}
  \UseMicrotypeSet[protrusion]{basicmath} % disable protrusion for tt fonts
}{}
\usepackage{setspace}
\makeatletter
\@ifundefined{KOMAClassName}{% if non-KOMA class
  \IfFileExists{parskip.sty}{%
    \usepackage{parskip}
  }{% else
    \setlength{\parindent}{0pt}
    \setlength{\parskip}{6pt plus 2pt minus 1pt}}
}{% if KOMA class
  \KOMAoptions{parskip=half}}
\makeatother
\setlength{\emergencystretch}{3em} % prevent overfull lines
\providecommand{\tightlist}{%
  \setlength{\itemsep}{0pt}\setlength{\parskip}{0pt}}
% v3 header: single column, tight reference list, clickable links.
% Loaded via pandoc -H. Fixes for the six issues Havid found in the v2 PDF are
% mostly in the markdown generator; this file only carries LaTeX-side settings.
\usepackage{graphicx}
\usepackage{float}
\usepackage{caption}

% Notation typesetting. Havid asked for consistency of format (2026-09-08),
% so units and species are expressed with the packages the field uses rather
% than a hand-rolled $^2$ regex table.
%
% VERIFIED before adoption: both compile under tectonic, which auto-downloads
% mhchem.sty and chemgreek.sty from its bundle -- no manual TeX Live install.
% Verified at PDF span level too: \qty{1.3e-3}{\centi\meter\squared\per\volt
% \per\second} raises every exponent AND every minus, and \ce{[PbI4]^2-} puts
% the subscript 4 BELOW the baseline with the charge ABOVE it in one token.
% A regex table cannot script one token in two directions.
%
% mhchem version=4 pinned: version 3 changed the charge syntax, so an
% unpinned load would silently alter every ion in the paper.
\usepackage{siunitx}
\usepackage[version=4]{mhchem}

% Journal-style units: upright glyphs, thin space between value and unit,
% and a real minus in exponents. detect-weight keeps units upright inside a
% bold heading instead of inheriting the surrounding bold.
\sisetup{
  detect-all,
  detect-weight        = true,
  range-units          = single,
  % reciprocal, NOT symbol. With per-mode=symbol the electron mobility
  % rendered as "cm2/(V s)", which is typographically correct but is not the
  % convention this literature uses; every perovskite paper writes
  % "cm2 V-1 s-1" with reciprocal powers. Verified at span level: reciprocal
  % raises both the 2 and each minus sign.
  per-mode             = reciprocal,
  exponent-product     = \times,
  output-decimal-marker = {.},
  group-digits         = none
}

% Figure captions: LaTeX owns the "Figure N." label. The v2 PDF rendered
% "Figure 4: Figure 4." because the caption text ALSO began with "Figure 4.",
% written into the markdown by the generator. v3 passes a bare caption and lets
% LaTeX number it, so the label appears exactly once.
\captionsetup[figure]{labelfont=bf, labelsep=period, font=small,
                      justification=justified, singlelinecheck=false,
                      skip=6pt}

% Tighten the numbered reference list produced by pandoc's default
% CSLReferences environment.
\usepackage{etoolbox}
\AtBeginEnvironment{CSLReferences}{%
  \small
  \setlength{\parskip}{0pt}%
  \setlength{\itemsep}{0pt}%
  \setlength{\topsep}{2pt}%
}

% No section numbering from LaTeX: the generator already writes "1. Introduction",
% "2. The Certified Frontier ...", so LaTeX numbering would double it.
\setcounter{secnumdepth}{0}
\usepackage{bookmark}
\IfFileExists{xurl.sty}{\usepackage{xurl}}{} % add URL line breaks if available
\urlstyle{same}
% fallback for those not using the hyperref driver hyperxmp:
\makeatletter
\@ifundefined{xmpquote}{\newcommand{\xmpquote}[1]{#1}}{}
\makeatother
\hypersetup{
  colorlinks=true,
  linkcolor={[HTML]{1A4E8A}},
  filecolor={Maroon},
  citecolor={Blue},
  urlcolor={[HTML]{1A4E8A}},
  pdfcreator={LaTeX via pandoc}}

\author{}
\date{}

\begin{document}

\setstretch{1.05}
\begin{center}
{\LARGE\bfseries Perovskite Photovoltaics in August 2026: Non-radiative recombination control and buried-interface treatment\par}
\vspace{1.1em}
{\large Havid Aqoma\textsuperscript{1,2,3,4,*}\par}
\vspace{0.7em}
\begin{minipage}{\textwidth}\centering\small
\textsuperscript{1}School of Energy and Chemical Engineering, Xiamen University Malaysia, Selangor Darul Ehsan 43900, Malaysia\\[0.25em]
\textsuperscript{2}Kelip-kelip! Center of Excellence for Light Enabling Technologies, School of Energy and Chemical Engineering, Xiamen University Malaysia, Selangor Darul Ehsan 43900, Malaysia\\[0.25em]
\textsuperscript{3}Centre of Excellence for Green and Advanced Technologies in Efficient Separations (GATES), School of Energy and Chemical Engineering, Xiamen University Malaysia, Selangor Darul Ehsan 43900, Malaysia\\[0.25em]
\textsuperscript{4}Gulei Innovation Institute, Xiamen University, Zhangzhou 363200, China\\[0.55em]
*Corresponding author. Email: havidaqoma.khoiruddin@xmu.edu.my\\[2.0em]
September 9, 2026
\end{minipage}
\end{center}
\vspace{0.8em}

\subsection{Abstract}\label{abstract}

Across a monthly literature corpus of 531 publications, 144 were
examined closely to evaluate defect passivation, contact chemistry,
phase stability, and operational durability in perovskite solar cells
(PSCs). Independent certification establishes the performance frontier
for monolithic tandem architectures at a power conversion efficiency
(PCE) of 33.66\%, driven by complementary bandgap absorption that
suppresses thermalization losses alongside optimized current matching.
Inverted single-junction configurations established a certified PCE
ceiling of 27.4\%, achieved primarily through tailored self-assembled
monolayers (SAMs) and molecular additives that passivate
undercoordinated lead defects and suppress interfacial molecular
clustering. Scaling these passivated contacts to expanded dimensions
incurs an area penalty governed by transparent conducting oxide sheet
resistance, interconnect scribe dead area, and spatial film
non-uniformity, yielding an accredited PCE of 17.35\% at the largest
reported aperture area of \qty{61.2}{\centi\meter}\(^2\). Operational
longevity evidence remains sparse and disconnected from efficiency
milestones. Only 5 investigations quantified operational lifetime, with
a maximum duration of \qty{1400}{\hour}, while only 2 studies adhered to
standardized International Summit on Organic Photovoltaic Stability
(ISOS) protocols. Degradation pathways identified across this cohort
include interfacial redox reactions between acidic phosphonic acid
anchoring groups and polar processing solvents, photo-induced halide
oxidation and ion migration, and residual lattice microstrain that
accelerates phase segregation under continuous operation. Reporting
practices reveal a pronounced discrepancy between self-reported claims
and certified performance, as 55.5\% of abstracts stated an efficiency
value, whereas only 9.7\% stated independent certification and only 18
articles reported verified values. Intermediate-phase modulation and
seed-layer crystallization mitigate interfacial voids and non-radiative
recombination at hole-selective contacts during initial deposition.
Published evidence fails to confirm that these strategies arrest
field-driven halide redistribution or contact degradation during
prolonged operation. Resolving these bottlenecks requires operando
chemical characterization under bias combined with standardized ISOS
testing on scalable apertures.

\textbf{Keywords:} perovskite photovoltaics; buried interface; defect
passivation; wide-bandgap perovskite; tandem solar cells; operational
stability

\subsection{1. Introduction}\label{introduction}

Perovskite photovoltaic research has moved beyond maximizing
unencapsulated power conversion efficiency (PCE) on sub-millimetre
laboratory apertures. The strongest mechanistic consensus demonstrates
that coordinating Lewis-base ligands and self-assembled monolayers
(SAMs) passivate non-radiative recombination centres at buried contacts
by coordinating undercoordinated lead sites. Yet the critical open
problem has shifted to operational lifetime, device area, and rigorous
verification. Arresting field-driven halide migration, phase
segregation, and contact corrosion during prolonged continuous operation
across scalable apertures remains unresolved, presenting barriers that
transient efficiency improvements cannot circumvent.

To evaluate these challenges, 531 indexed works met the scope gate in
August 2026 and 144 were read closely.

The resulting evidence base is structured across core physical domains:
defect tolerance, passivation and ion migration, buried interfaces and
contact chemistry, composition, phase control and crystallisation,
wide-bandgap absorbers and tandem integration, as well as operational
stability and degradation pathways.

Throughout each topic, this review distinguishes what the month's
literature verified from what it asserted. Theoretical passivation
mechanisms and charge-transport claims are scrutinized directly against
empirical device behavior, separating documented operational
survivability from unverified assertions based on uncertified,
small-area cells subjected to brief, non-standardized aging routines.

\subsection{2. The Certified Frontier in August
2026}\label{the-certified-frontier-in-august-2026}

Independent certification establishes that spectral splitting across
complementary bandgaps yields the highest verified efficiencies.
Suppressed thermalization and optimized current matching drive this
outcome in two-terminal (2T) perovskite tandems. Verified 2T power
conversion efficiency (PCE) benchmarks reached a ceiling of 33.66\%,
followed by certified reports at 32.56\%, 30.77\%, 29.7\%, and 27.84\%
{[}\href{https://doi.org/10.1038/s41467-026-76266-0}{1},\href{https://doi.org/10.1002/anie.5241611}{2},\href{https://doi.org/10.1038/s41467-026-76713-y}{3},\href{https://doi.org/10.1016/j.joule.2026.102645}{4},\href{https://doi.org/10.1038/s41467-026-76786-9}{5}{]}.
Single-junction inverted (p-i-n) devices, constrained by non-radiative
interfacial recombination and extraction barriers at selective contacts,
established a narrower certified ceiling between 26.77\% and 27.4\%, led
by verified values of 27.4\%, 27.2\%, and 26.77\%
{[}\href{https://doi.org/10.1126/science.aeg8415}{6},\href{https://doi.org/10.1038/s41467-026-76497-1}{7},\href{https://doi.org/10.1002/anie.3095986}{8}{]}.
Despite these verified milestones, rigorous certification remains rare
across the literature: only 18 of 144 closely read papers reported an
independently certified efficiency, allowing uncertified self-reported
champion values to run ahead of verifiable baselines.

Translating these contact interfaces across larger dimensions exposes
severe transport bottlenecks governed by transparent electrode sheet
resistance, scribe line geometric dead area, and spatial film
non-uniformity. At expanded dimensions, minimodule architectures
maintained reported efficiencies of 27.7\% on
\qty{25}{\centi\meter\squared} and 27.01\% on
\qty{58.51}{\centi\meter\squared}
{[}\href{https://doi.org/10.1038/s43246-026-01321-8}{9},\href{https://doi.org/10.1002/adma.74469}{10}{]}.
However, larger apertures revealed significant performance losses,
dropping to 23.3\% on \qty{57.3}{\centi\meter\squared}, 23.14\% on
\qty{29.7}{\centi\meter\squared}, and 17.35\% on
\qty{61.2}{\centi\meter\squared}
{[}\href{https://doi.org/10.1002/adma.74629}{11},\href{https://doi.org/10.1002/advs.76950}{12},\href{https://doi.org/10.1002/smll.75068}{13}{]}.
A champion value recorded on a sub-\qty{0.1}{\centi\meter\squared}
laboratory cell isolates microscopic interfacial transport under
near-ideal current extraction, whereas an independently certified value
measured across a module aperture gauges lateral conduction, spatial
defect distributions, and coating reproducibility. These two metrics
represent entirely distinct categories of physical evidence, meaning
small-cell records cannot demonstrate module viability.

\begin{figure}[htbp]
\centering
\includegraphics[width=0.95\textwidth]{fig/F1_certified_frontier.pdf}
\caption{The certified frontier in August 2026. Filled markers are independently certified efficiencies, open markers are self-reported champion values, grouped by device family. Horizontal bars mark the highest certified value in each family.}
\end{figure}

\subsection{3. Defect Tolerance, Passivation and Ion
Migration}\label{defect-tolerance-passivation-and-ion-migration}

Defect formation and ion migration in metal halide perovskites are
coupled through mobile halide vacancies that trigger operational
degradation. Continuous illumination accelerates iodide oxidation to
molecular iodine, generating iodine vacancies that expose
undercoordinated lead (\ce{Pb^2+}) sites
{[}\href{https://doi.org/10.1021/acsenergylett.6c02033}{14}{]}, whereas
nonadiabatic molecular dynamics shows that iodine vacancies drive
hydrogen vacancy formation into stable VI·VH defect complexes whose
migration is immobilized by rigid Pb-C coordinate bonding in VI·VH-C
configurations {[}\href{https://doi.org/10.1021/acsnano.6c10992}{15}{]}.
Kinetic and compositional regulation confirms that mobile halide
management is prerequisite to suppressing ion drift
{[}\href{https://doi.org/10.1002/adfm.77855}{16},\href{https://doi.org/10.1002/aenm.71462}{17}{]}.
Exploiting the kinetic accessibility of bromide (Br-) diffusion, an
ordered anion-cation modification (OACM) strategy repairs anionic
vacancies before cation passivation, suppressing phase segregation to
achieve a \qty{1.33}{\volt} open-circuit voltage (Voc) and 19.21\%
stabilized efficiency in \qty{1.83}{\electronvolt} wide-bandgap (WBG)
top cells (champion 26.2\% tandem efficiency)
{[}\href{https://doi.org/10.1002/adfm.77855}{16}{]}. Selenocyanate
pseudohalide alloying similarly retards ion drift and homogenizes
formamidinium (FA+) distributions throughout \qty{1.66}{\electronvolt}
absorbers {[}\href{https://doi.org/10.1002/aenm.71462}{17}{]}.
Structural barriers provide complementary confinement. Scaffold-directed
polyfluorinated 7F-EA-HI forms an interwoven tilted
one-dimensional/three-dimensional (1D/3D) heterointerface that blocks
defect propagation and ion migration in inverted (p-i-n) devices
reaching an uncertified 24.94\% power conversion efficiency (PCE)
{[}\href{https://doi.org/10.1002/anie.8256065}{18}{]}. Self-driven
halide exchange with bromide-rich nanocrystals passivates iodide
vacancies via a surface-graded heterojunction, though restricted to an
uncertified 22.19\% champion PCE in an unreviewed preprint
{[}\href{https://doi.org/10.26434/chemrxiv.15007814/v1}{19}{]}.

Beyond mobile halides, eliminating deep non-radiative traps demands
coordinated passivation of undercoordinated metal cations and organic
cation vacancies
{[}\href{https://doi.org/10.1021/acsenergylett.6c02033}{14},\href{https://doi.org/10.1002/adfm.78067}{20},\href{https://doi.org/10.1021/acsami.6c10456}{21},\href{https://doi.org/10.1021/acs.jpclett.6c02253}{22}{]}.
Bulk phenylthiourea coordinates \ce{Pb^2+} to retard ion migration and
suppress defect formation under illumination
{[}\href{https://doi.org/10.1021/acsenergylett.6c02033}{14}{]}. Ammonium
alginate suppresses FA+ deprotonation while passivating undercoordinated
\ce{Pb^2+}, raising fill factor (FF) from 80.2\% to 85.1\% in an
uncertified 26.12\% champion cell
{[}\href{https://doi.org/10.1002/adfm.78067}{20}{]}. Dual-site
coordination similarly targets opposing charges: sodium cyclamate
coordinates \ce{Pb^2+} via sulfonic acid groups and binds iodide
vacancies through amino hydrogen bonds
{[}\href{https://doi.org/10.1021/acsami.6c10456}{21}{]}, while
multidentate sulfamethoxazole coordinates \ce{Pb^2+}, organic-cation,
and halide vacancies, lifting normal (n-i-p) efficiency from 18.01\% to
an uncertified 21.05\% champion (20.31\% average)
{[}\href{https://doi.org/10.1021/acs.jpclett.6c02253}{22}{]}. Coupled
hydrogen-bonding and electrostatic networks from the ionic liquid BSPH
coordinate FA+, halides, and \ce{Pb^2+}, suppressing trap density and
ion drift to yield a 23.72\% p-i-n PCE
{[}\href{https://doi.org/10.1002/solr.70461}{23}{]}. In
hole-conductor-free carbon-electrode mesoscopic cells, DMSC passivates
both Pb and I species to reach a 22.5\% champion PCE
{[}\href{https://doi.org/10.1002/anie.5667084}{24}{]}, while
multi-pyridyl spirobifluorenes coordinate defects during crystallization
to yield 20.98\% PCE versus 18.67\% in controls
{[}\href{https://doi.org/10.1021/acsami.6c13903}{25}{]}. In tin-based
absorbers prone to \ce{Sn^2+} oxidation, triphenylphosphine oxide
coordinates \ce{Sn^2+} as a hard Lewis base, yielding 14.13\% PCE and a
T80 of \qty{720}{\hour} under unlabelled testing conditions
{[}\href{https://doi.org/10.1002/solr.70454}{26}{]}.

A third frontier links defect formation directly to lattice mechanics
and buried transport contacts
{[}\href{https://doi.org/10.1093/nsr/nwag484}{27},\href{https://doi.org/10.1002/aenm.71497}{28},\href{https://doi.org/10.1002/smll.75360}{29}{]}.
Solid-solution lattice hardening by partial lanthanum substitution of Pb
raises the Young modulus from 27.96 to 35.71 GPa, demonstrating that
rigidifying the inorganic cage suppresses defect formation to yield a
26.54\% champion PCE and a certified steady-state efficiency of 26.03\%
{[}\href{https://doi.org/10.1093/nsr/nwag484}{27}{]}. Dynamic strain
homogenization via in situ supramolecular networks modulates \ce{Sn^2+}
coordination in mixed tin-lead absorbers, suppressing segregation to
retain 84.7\% initial efficiency after \qty{450}{\hour} maximum power
point tracking (MPPT) under one sun and 92.1\% after 50 thermal cycles
(24.1\% champion PCE)
{[}\href{https://doi.org/10.1002/aenm.71497}{28}{]}. At buried contacts,
Ace-K bidentately chelates uncoordinated \ce{Sn^4+} and oxygen vacancies
on tin dioxide (\ce{SnO2}) while outcompeting dimethyl sulfoxide (DMSO)
for lead iodide (\ce{PbI2}), eliminating voids to retain 92.4\%
efficiency after \qty{1000}{\hour} MPPT on a
\qty{0.0729}{\centi\meter\squared} cell (26.17\% champion PCE)
{[}\href{https://doi.org/10.1002/smll.75360}{29}{]}. Interfacial
molecular bridging with 2-formylbenzenesulfonic acid sodium salt
connects self-assembled monolayers (SAM) to perovskite, delivering an
uncertified 26.21\% champion PCE with an 86.15\% FF
{[}\href{https://doi.org/10.1007/s40820-026-02278-6}{30}{]}. At module
scale, potassium dihydrogen phosphate disaggregates \ce{SnO2} colloids
and passivates buried contacts, yielding 23.16\% PCE on
\qty{12}{\centi\meter\squared} modules with 8\% scaling loss from a
25.34\% champion cell
{[}\href{https://doi.org/10.1002/aenm.71419}{31}{]}.

\begin{figure}[htbp]
\centering
\includegraphics[width=0.95\textwidth]{fig/F3_effort_map.pdf}
\caption{Where the month's mechanistic effort was directed. Depth-tier papers by mechanism axis, stacked by device family.}
\end{figure}

\subsection{4. Buried Interfaces and Contact
Chemistry}\label{buried-interfaces-and-contact-chemistry}

At buried hole-selective contacts in inverted
positive-intrinsic-negative (p-i-n) and tandem perovskite solar cells
(PSCs), disordered molecular packing and aggregation in self-assembled
monolayers (SAMs) cause severe pinhole shunts and energy misalignment.
Suppressing self-assembled monolayer (SAM) clustering while keeping
anchoring sites accessible resolves these losses: non-anchoring
additives like TTMB secure open binding sites to reach a certified power
conversion efficiency (PCE) of 33.66\% (33.66\% stabilised) in
two-terminal (2T) perovskite/silicon tandems
{[}\href{https://doi.org/10.1038/s41467-026-76266-0}{1}{]}.
Hydrogen-bond networks via glycinamide hydrochloride in
{[}2-(3,6-dimethoxy-9H-carbazol-9-yl)ethyl{]}phosphonic acid (MeO-4PACz)
yield 22.38\% champion PCE in pure-iodide wide-bandgap cells
{[}\href{https://doi.org/10.1016/j.decarb.2026.100190}{32}{]}.
DMPA-based multilayers achieve a certified 27.2\% PCE (27.59\% champion)
with 94.5\% retention after 1600 hours under continuous 1 Sun
illumination at 65 °C
{[}\href{https://doi.org/10.1038/s41467-026-76497-1}{7}{]}. Interfacial
dipole engineering further optimizes band offsets. Depositing bismuth
oxychloride (BiOCl) onto SAMs induces dipole flipping and
chlorine-iodine exchange, yielding a certified 27.84\% 2T tandem PCE
(28.14\% champion, 27.86\% stabilised), although single-junction
efficiencies of 27.25\% (\qty{1.55}{\electronvolt}) and 20.24\%
(\qty{1.85}{\electronvolt}) remain uncertified on small
\qty{0.09}{\centi\meter\squared} aperture areas
{[}\href{https://doi.org/10.1038/s41467-026-76786-9}{5}{]}. This aligns
with co-assembling 4-ASA with Me-4PACz on nickel oxide (\ce{NiO_x}) to
yield a 20.36\% four-terminal (4T) tandem PCE (17.82\% scanned single
junction), and dipole tuning via m-FBC raising PCE from 23.05\% to an
uncertified 25.03\%, though stability was tested only under 500 hours of
uncharacterized storage
{[}\href{https://doi.org/10.1038/s41467-026-76847-z}{33},\href{https://doi.org/10.1021/acsami.6c11467}{34}{]}.
Asymmetric heteroatom conjugation also regulates contact energetics:
S-AS on \ce{NiO_x} achieves 25.70\% steady-state output (26.49\%
champion) and LW-1 on fluorine-doped tin oxide (FTO) reaches 26.0\%
stabilised efficiency (26.38\% champion)
{[}\href{https://doi.org/10.1002/adfm.77712}{35},\href{https://doi.org/10.1002/aenm.71482}{36}{]}.

Beyond electrostatic alignment, buried contact molecules template
crystallization, relieve lattice strain, and passivate undercoordinated
lead (Pb2⁺) defects. Multi-arm donor-acceptor dipoles (T4IA) coordinate
Pb2⁺ through C=N and methoxy groups to relieve residual strain (26.79\%
champion PCE), whereas non-covalent templating between precursor
additives (TPANI) and structurally matched triphenylamine bisphosphonic
acid SAMs on indium tin oxide (ITO) reduces intergranular groove depth
and reinforces adhesion (26.58\% champion)
{[}\href{https://doi.org/10.1002/anie.3716107}{37},\href{https://doi.org/10.1002/anie.6932848}{38}{]}.
Interfacial Pb-S bonding provides a parallel passivation route.
Thienyl-based SAM additives template preferential (111) crystallization
and segregate at grain boundaries to secure a certified 32.56\% (33.0\%
champion) 2T tandem PCE, mirroring solution-phase sulfidation with
dibutyl sulfide that induces an interfacial p-to-n transition to yield
an uncertified 26.28\% PCE
{[}\href{https://doi.org/10.1002/anie.5241611}{2},\href{https://doi.org/10.1021/acsami.6c07884}{39}{]}.

Contact degradation also stems from chemical reactivity across metal
oxide substrates. On \ce{NiO_x} hole transport layers (HTLs), removing
deleterious surface hydroxyls with ammonium thiocyanate (\ce{NH4SCN})
lowers transfer barriers and lattice strain (21.41\% champion)
{[}\href{https://doi.org/10.1021/acsaem.6c01708}{40}{]}. Decoupling SAM
anchoring from charge extraction through composite \ce{NiO_x} stacks
(sputtered base plus solution-deposited hydroxyl overlayer) permits
scalable processing, achieving 24.86\% PCE across
\qty{12.4}{\centi\meter\squared}
{[}\href{https://doi.org/10.1002/smll.75056}{41}{]}. Enamine SAM anchor
engineering is consistent with this, achieving 23.14\% PCE with 83.11\%
fill factor (FF) on a \qty{29.7}{\centi\meter\squared} mini-module
(\textgreater25.5\% on small cells)
{[}\href{https://doi.org/10.1002/advs.76950}{12}{]}. At tin oxide
(\ce{SnO2}) electron transport layers (ETLs) in regular
negative-intrinsic-positive (n-i-p) devices, native surface hydroxyls
and oxygen vacancies provoke interfacial void formation and
microcracking. Converting hydroxyls into covalently anchored acetate via
ethanol vapor reconstruction suppresses both (26.19\% champion)
{[}\href{https://doi.org/10.1002/smll.75122}{42}{]}. Subliming pyruvic
acid at 145 °C purges defect sites instead, reaching 26.17\% PCE on
\qty{1}{\centi\meter\squared} with an 80\% operational lifetime (T80) of
900 hours {[}\href{https://doi.org/10.1002/aenm.71514}{43}{]}. However,
interfacial engineering reveals a trade-off between extraction and
mechanical integrity: although fullerene SAMs (C60-PA) on \ce{SnO2}
raise initial tracked efficiencies from 18.6\% to a 20.0\% stabilised
value, performance degrades under illumination at elevated temperatures
because fullerenes weaken adhesion at the oxide interface
{[}\href{https://doi.org/10.1021/acsenergylett.6c01485}{44}{]}.

\subsection{5. Composition, Phase Control and
Crystallisation}\label{composition-phase-control-and-crystallisation}

Inhomogeneous nucleation triggers point defects, phase segregation, and
structural disorder in metal halide perovskite absorbers, requiring
precise crystallization modulation
{[}\href{https://doi.org/10.1002/adma.74388}{45}{]}. Controlling
intermediate solvent-solute coordination dictates lattice orientation
and residual stress. Destabilizing solvent-\ce{PbI2} adducts through
cooperative interactions between sodium hydroxyquinoline sulfonate and
sodium aminonaphthalene sulfonate drives out-of-plane {[}100{]} growth
while converting tensile strain into compressive strain. This yields an
uncertified champion power conversion efficiency (PCE) of 26.21\%
{[}\href{https://doi.org/10.1002/aenm.71474}{46}{]}. Similarly,
dual-mode coordination with 4-chlorophenethylammonium (4Cl-PA) generates
a macroporous \ce{PbI2} template that facilitates organic cation
infiltration. Carrier lifetime extends to \qty{161}{\nano\second} by
time-resolved photoluminescence (TRPL) in uncertified 22.73\% efficient
devices {[}\href{https://doi.org/10.1016/j.jcis.2026.141449}{47}{]}. In
tin perovskites, intermediate-mediated stepwise crystallization during a
pre-annealing static stage with fluorinated phenethylammonium alters the
\ce{[SnI6]^4-} electronic landscape to suppress \ce{Sn^2+} oxidation,
reaching a 9.29\% PCE
{[}\href{https://doi.org/10.1021/acsami.6c11057}{48}{]}. However,
chemical coordination efficacy depends on the processing solvent.
\ce{SnF2} suppresses \ce{Sn^2+} oxidation in dimethylformamide:dimethyl
sulfoxide (DMF:DMSO). It offers limited antioxidant protection in
DMSO-free DMF:2-methylaminoethanol (DMF:2MAP), where formamidinium tin
iodide (FASnI3) films remain inherently \ce{Sn^2+}-rich
{[}\href{https://doi.org/10.1021/acsaem.6c02063}{49}{]}.

Phase stability is equally bounded by thermodynamic competition and
local defect equilibria
{[}\href{https://doi.org/10.1002/adma.74538}{50},\href{https://doi.org/10.1002/aenm.71396}{51},\href{https://doi.org/10.1002/advs.76791}{52}{]}.
In cesium-based systems, excess \ce{PbI2} suppresses Ruddlesden-Popper
antiphase boundaries (RP-APBs) but triggers edge-sharing \ce{[PbI6]^4-}
motifs that nucleate the non-photoactive delta phase. Dimethylammonium
(DMA+) phase engineering circumvents this breakdown by forming
beta-(DMA,Cs)PbI3 and intrinsically eliminating RP-APBs, achieving an
uncertified 20.23\% PCE
{[}\href{https://doi.org/10.1002/adma.74538}{50}{]}. Similarly,
incorporating 0.8 mol\% fluorinated methylammonium iodide (FMAI)
stabilizes the perovskite lattice against thermal degradation in a
25.67\% negative-intrinsic-positive (n-i-p) device. The reported
\qty{2000}{\hour} operational stability lacks detailed testing protocols
{[}\href{https://doi.org/10.1002/aenm.71396}{51}{]}. In low-dimensional
heterostructures, organic cation evaporation during annealing drives a
dimensional shift from low- to high-n quantum wells. Precursor
octylammonium bromide (OABr) concentration predetermines the initial n =
1 phase distribution in 19.29\% wide-bandgap cells
{[}\href{https://doi.org/10.1002/advs.76791}{52}{]}. Replacing
insulating organic spacers with semiconducting diketopyrrolopyrrole
cations enables bidirectional charge transfer with inorganic
\ce{[PbI4]^2-} slabs. Electron mobilities reach up to
\qty{1.3e-3}{\centi\meter\squared\per\volt\per\second}, although
functional solar cell metrics were omitted
{[}\href{https://doi.org/10.1002/adma.74681}{53}{]}.

In lead-free multinary lattices, accessing target crystal phases
requires overcoming severe kinetic and thermodynamic barriers
{[}\href{https://doi.org/10.1016/j.solener.2026.114961}{54},\href{https://doi.org/10.26434/chemrxiv.15008057/v1}{55},\href{https://doi.org/10.1038/s43246-026-01261-3}{56}{]}.
Sequentially evaporating methylbismuth iodide (MBI) bilayers at
substrate temperatures as low as 40 °C requires post-deposition
annealing above 200 °C to yield a monoclinic C2/c phase with (111)
texture. Champion PCE nonetheless reaches only 0.53\%
{[}\href{https://doi.org/10.1016/j.solener.2026.114961}{54}{]}. In
double perovskites, solid-state mechanosynthesis reveals strict ionic
size selectivity. Liquid assistance drives \ce{Rb+} into the A-site and
\ce{Na+} into the B-site of \ce{Cs2AgInCl6}
{[}\href{https://doi.org/10.26434/chemrxiv.15008057/v1}{55}{]}.
Conversely, solvent-free hydro-solution precipitation directly produces
cubic Fm-3m \ce{Cs2AgBiBr6} without high-temperature annealing
{[}\href{https://doi.org/10.1038/s43246-026-01261-3}{56}{]}. This
indicates that solution-phase ionic transport bypasses the kinetic
barriers inherent to vapor-phase and solid-state transformations.

\subsection{6. Wide-Bandgap Absorbers and Tandem
Integration}\label{wide-bandgap-absorbers-and-tandem-integration}

Non-conformal coverage and dewetting of hole-selective self-assembled
monolayers (SAMs) on textured silicon substrates generate interfacial
voids and unreacted lead iodide (\ce{PbI2}), driving non-radiative
recombination in two-terminal (2T) perovskite/silicon tandems
{[}\href{https://doi.org/10.1016/j.joule.2026.102645}{4},\href{https://doi.org/10.1002/aenm.71512}{57},\href{https://doi.org/10.1021/acsenergylett.6c01966}{58}{]}.
Resolving this buried morphological mismatch requires tailored
nucleation and crystallization. A CsCl seed layer facilitates organic
precursor incorporation across MeO-2PACz to suppress interfacial
\ce{PbI2}, yielding a certified 29.7\% power conversion efficiency (PCE)
{[}\href{https://doi.org/10.1016/j.joule.2026.102645}{4}{]}.
Solvent-free sequential evaporation capitalizes on accelerated lead
halide conversion on textured silicon to deliver a 27.0\% stabilized PCE
{[}\href{https://doi.org/10.1016/j.joule.2026.102646}{59}{]}.
Post-deposition methylammonium thiocyanate (MASCN) healing similarly
repairs valley-confinement defects, driving secondary crystal growth
into full-thickness columnar grains. That device achieved a certified
30.77\% PCE on a \qty{1.164}{\centi\meter\squared} aperture and retained
initial performance after 3400 hours of continuous ambient maximum power
point tracking (MPPT)
{[}\href{https://doi.org/10.1038/s41467-026-76713-y}{3}{]}. In the
wide-bandgap (WBG) top absorber, phase segregation and halide migration
govern bulk degradation. Elevating the ion migration energy barrier
improves local structural stability, supporting an uncertified champion
PCE of 33.02\% {[}\href{https://doi.org/10.1002/adfm.77637}{60}{]}.
Two-dimensional (2D) (APD)PbI4 seed templates impart regular compressive
strain into the three-dimensional (3D) lattice, suppressing phase
separation for a 29.47\% PCE
{[}\href{https://doi.org/10.1021/acsnano.6c07355}{61}{]}. Contact
durability under subsequent processing presents a further constraint.
Sacrificial HTFNA coordination shields SAMs against atomic layer
deposition (ALD) damage, yet this certified 26.56\% efficiency is
restricted to a \qty{0.062}{\centi\meter\squared} cell
{[}\href{https://doi.org/10.1002/adma.74461}{62}{]}. The field is
shifting toward resilient inorganic contacts like nickel oxide with N719
dye, which preserve 96.4\% PCE after one year of outdoor MPPT
{[}\href{https://doi.org/10.1021/acsenergylett.6c01363}{63}{]}.
Translating these interfaces to large-area modules escalates multi-scale
defect densities
{[}\href{https://doi.org/10.1002/adfm.202532038}{64}{]}, requiring
kinetic crystallization control. Polyfluoroarene additives elevate
effective nucleation barriers to yield a 27.01\% module PCE across
\qty{58.51}{\centi\meter\squared} with \qty{55}{\milli\volt}
non-radiative voltage loss
{[}\href{https://doi.org/10.1002/adma.74469}{10}{]}. Static antisolvent
engineering matches this, achieving 27.7\% on
\qty{25}{\centi\meter\squared}
{[}\href{https://doi.org/10.1038/s43246-026-01321-8}{9}{]}.

\begin{figure}[htbp]
\centering
\includegraphics[width=0.95\textwidth]{fig/F2_area_penalty.pdf}
\caption{Efficiency against reported active or aperture area. The shaded band marks the sub-0.3 cm$^2$ laboratory-cell regime in which most champion values are measured.}
\end{figure}

\subsection{7. Operational Stability and Degradation
Pathways}\label{operational-stability-and-degradation-pathways}

Interfacial chemical and photochemical degradation at selective contacts
constitutes the primary failure mechanism addressed this month. Inverted
perovskite solar cells (PSCs) suffer from coupled redox degradation.
Acidity in carbazole-based phosphonic acid self-assembled monolayers
(SAMs) activates dimethyl sulfoxide (DMSO)-mediated iodide redox
reactions during ink deposition
{[}\href{https://doi.org/10.1126/science.aeg8415}{6}{]}. Photo-oxidation
of formamidinium iodide (FAI) yields iodine radicals that catalyze
phenyl-C61-butyric acid methyl ester (PCBM) dimerization alongside FAI
deprotonation {[}\href{https://doi.org/10.1002/anie.4221605}{65}{]}.
Cross-linking self-assembled monolayer (SAM) headers in situ
{[}\href{https://doi.org/10.1126/sciadv.aec3096}{66}{]}, co-assembling
binary monolayers
{[}\href{https://doi.org/10.1038/s41467-026-76848-y}{67}{]}, or
balancing dipole configurations
{[}\href{https://doi.org/10.1002/adfm.77942}{68},\href{https://doi.org/10.1002/adma.74438}{69}{]}
suppresses interfacial defect generation. Yet performance claims
frequently exhibit severe certification discrepancies or rely on tiny
device areas, such as a 27.06\% uncertified champion dropping to 20.11\%
certified on \qty{0.09}{\centi\meter\squared}
{[}\href{https://doi.org/10.1002/adma.74438}{69}{]} and a certified
27.4\% power conversion efficiency (PCE) tested on
\qty{0.06}{\centi\meter\squared}
{[}\href{https://doi.org/10.1126/science.aeg8415}{6}{]}.

In the bulk absorber, unrelaxed microstrain and defect formation
accelerate phase degradation. Three routes relieve lattice strain: 2
mol\% sodium (\ce{Na+}) interstitials
{[}\href{https://doi.org/10.1002/anie.2425823}{70}{]}, two-dimensional
(2D) intermediate phase templating
{[}\href{https://doi.org/10.1038/s41467-026-76425-3}{71}{]}, and
vertical lattice misalignment suppression
{[}\href{https://doi.org/10.1002/adfm.77729}{72}{]}. Each increases
defect formation energies for lead and iodide vacancies and prevents
nonphotoactive phase transitions
{[}\href{https://doi.org/10.1002/anie.2425823}{70},\href{https://doi.org/10.1038/s41467-026-76425-3}{71},\href{https://doi.org/10.1002/adfm.77729}{72}{]}.
At rear contacts, a \qty{5}{\nano\meter} chromium (Cr) barrier beneath
non-gold electrodes curtails atomic interdiffusion
{[}\href{https://doi.org/10.1002/aenm.71372}{73}{]}. Coordination
modifiers prevent direct perovskite contact with metal oxides
{[}\href{https://doi.org/10.1002/adfm.77805}{74}{]}.

Standardized operational metrics remain conspicuously scarce. Across 144
examined papers, only 5 reported a time to 80\% of initial efficiency
(T80). Merely 2 adhered to International Summit on Organic Photovoltaic
Stability (ISOS) protocols, specifically ISOS-L-1
{[}\href{https://doi.org/10.1038/s41467-026-76848-y}{67}{]} and ISOS-L-2
{[}\href{https://doi.org/10.1002/anie.4221605}{65}{]}. Consequently,
isolated metrics such as an \qty{800}{\hour} T80
{[}\href{https://doi.org/10.1002/anie.2425823}{70}{]} or an
\qty{864}{\hour} T80 on \qty{0.08}{\centi\meter\squared}
{[}\href{https://doi.org/10.1002/adfm.77805}{74}{]} reflect
unstandardized evaluations rather than verified operational longevity.

\begin{figure}[htbp]
\centering
\includegraphics[width=0.95\textwidth]{fig/F4_stability_evidence.pdf}
\caption{The month's operational-stability evidence in full. Left, every reported time to 80\% of initial performance. Right, the stability protocol labels used.}
\end{figure}

\subsection{8. Research Gaps and
Outlook}\label{research-gaps-and-outlook}

Controlling wide-bandgap (WBG) perovskite crystallization across
textured silicon substrates represents the most supported mechanism this
month, advancing certified two-terminal (2T) power conversion efficiency
(PCE) to 33.66\%
{[}\href{https://doi.org/10.1038/s41467-026-76266-0}{1}{]} and 30.77\%
{[}\href{https://doi.org/10.1038/s41467-026-76713-y}{3}{]}. However, the
assembled evidence leaves critical physical questions unestablished
across four distinct areas.

First, although seed layers and secondary grain growth suppress phase
segregation during deposition
{[}\href{https://doi.org/10.1038/s41467-026-76713-y}{3},\href{https://doi.org/10.1016/j.joule.2026.102645}{4}{]},
current data do not demonstrate that columnar grains prevent
electric-field-driven halide redistribution under sustained bias.
Operando depth-profiling time-of-flight secondary ion mass spectrometry
(ToF-SIMS) under operational voltage bias is required to determine
whether seeded interfaces permanently impede ion migration or merely
delay initial demixing. Second, non-anchoring additives reduce
aggregation in self-assembled monolayers (SAMs)
{[}\href{https://doi.org/10.1038/s41467-026-76266-0}{1}{]}, yet the
literature fails to verify whether these modified contacts survive wet
chemical processing. Acidic phosphonic acid SAMs trigger dimethyl
sulfoxide (DMSO)-mediated iodide redox reactions during ink coating
{[}\href{https://doi.org/10.1126/science.aeg8415}{6}{]}, leaving
uncertified champion efficiencies vulnerable to severe certification
discrepancies {[}\href{https://doi.org/10.1002/adma.74438}{69}{]}. In
situ grazing-incidence wide-angle X-ray scattering (GIWAXS) coupled with
interface-sensitive photoelectron spectroscopy during solution wetting
must establish whether molecular additives protect buried contacts from
chemical degradation. Third, device scalability remains disconnected
from small-area interfacial gains. Scale-up was addressed by 27 of 531
works, with the largest reported aperture reaching
\qty{61.2}{\centi\meter\squared}
{[}\href{https://doi.org/10.1002/smll.75068}{13}{]}. Reported data still
do not decouple the physical origins of area-dependent efficiency
penalties {[}\href{https://doi.org/10.1038/s41467-026-76425-3}{71}{]}.
Calibrated absolute electroluminescence imaging resolving spatial
ideality factors across square-decimeter modules is necessary to
separate interconnect sheet resistance from spatial non-radiative
recombination. Fourth, operational longevity remains unverified across
the field. Only 5 papers reported a time to 80\% of initial performance
(T80) lifetime, with the longest reaching \qty{1400}{\hour}
{[}\href{https://doi.org/10.1002/anie.6548693}{75}{]}. Only 2 papers
named an International Summit on Organic Photovoltaic Stability (ISOS)
protocol
{[}\href{https://doi.org/10.1002/anie.4221605}{65},\href{https://doi.org/10.1038/s41467-026-76848-y}{67}{]}.
Standardized ISOS testing combining elevated thermal stress with
continuous maximum power point tracking (MPPT) under full illumination
is needed to isolate contact delamination from intrinsic lattice decay.

The certified reporting share was 8.3 in 2026-07 and 9.7 in 2026-08, yet
only 18 of 144 closely read papers reported any certified value.
Requiring standardized ISOS protocols alongside certified aperture-area
efficiencies would make next month's literature directly comparable
across laboratories.

\subsection{Acknowledgements}\label{acknowledgements}

The authors would like to acknowledge the financial support from Xiamen
University Malaysia through the Xiamen University Malaysia Research Fund
(XMUMRF/2022-C10/IENG/0046), and from National Natural Science
Foundation of China (Ref. No.: 52503390).

\subsection{Declaration of Competing
Interest}\label{declaration-of-competing-interest}

The authors declare no conflict of interest.

\subsection{Data Availability
Statement}\label{data-availability-statement}

The corpus table for the month, the per-paper extraction records with
their verbatim source quotations, the statistics file underlying every
number in the text, and the gate report are provided as Supplementary
Information. Abstract text and full texts are not redistributed.

\subsection{AI Usage Declaration}\label{ai-usage-declaration}

This manuscript was prepared as part of our Research Project for making
a reliable Autonomous Researcher Agent. The agent executed the monthly
review pipeline without human intervention: it harvested the indexed
literature for the month through the OpenAlex, Semantic Scholar,
Crossref and arXiv interfaces, applied a deterministic topical filter,
ranked and selected the papers read closely, extracted every
quantitative claim into a structured record binding each number to a
verbatim quotation from that paper's own abstract, computed the
statistics and figures from those records, and drafted this text from
the resulting evidence briefs. Literature harvesting, selection,
statistics, figures and gate checking are deterministic scripts rather
than model output. Structured extraction used qwen3.8-flash through the
opencode interface; the prose was written by gemini-3.8-flash-high
through the agy interface; independent review used Opus 5 at high
reasoning effort. No figure is a model-generated image: every panel is
plotted directly from the extracted records. Every quantitative
statement in the main text resolves to the machine-generated statistics
file shipped as Supplementary Information, and every citation resolves
to a record carrying its source quotation. The author reviewed and
edited the content, verified the selection, and takes full
responsibility for the content of the published article.

\subsection{References}\label{references}

\begin{enumerate}
\def\labelenumi{\arabic{enumi}.}
\tightlist
\item
  Zeng, L. et al.~Non-competitive additives regulate molecular assembly
  for compact monolayers in perovskite/silicon tandems. \emph{Nature
  Communications} (2026).
  \url{https://doi.org/10.1038/s41467-026-76266-0}
\item
  Zhang, D. et al.~Thienyl Self‐Assembled Molecule Enabled Stable
  Perovskite/TOPCon Silicon Tandem. \emph{Angewandte Chemie
  International Edition} (2026).
  \url{https://doi.org/10.1002/anie.5241611}
\item
  Kong, W. et al.~Healing intervention for improving the efficiency and
  stability of tandem devices on industrial textured silicon.
  \emph{Nature Communications} (2026).
  \url{https://doi.org/10.1038/s41467-026-76713-y}
\item
  Škorjanc, V. et al.~CsCl seed layer homogenizes co-evaporated
  perovskite growth for high-efficiency, fully textured
  perovskite-silicon tandem solar cells. \emph{Joule} (2026).
  \url{https://doi.org/10.1016/j.joule.2026.102645}
\item
  Xie, Z. et al.~Stacking-order-driven interfacial dipole flip and
  chlorine-iodine exchange for high-performance perovskite-organic
  tandem solar cells. \emph{Nature Communications} (2026).
  \url{https://doi.org/10.1038/s41467-026-76786-9}
\item
  Liang, Z. et al.~Redirecting wet--interfacial redox pathways for
  efficient inverted perovskite solar cells. \emph{Science} (2026).
  \url{https://doi.org/10.1126/science.aeg8415}
\item
  Pan, Y. et al.~Self-Assembled Multilayers Reduce Interfacial Energy
  Loss in Perovskite Solar Cells. \emph{Nature Communications} (2026).
  \url{https://doi.org/10.1038/s41467-026-76497-1}
\item
  Feng, Y. et al.~Ga Doping Enables (110) Facet Control and Surface
  Passivation of SnO 2 for Efficient and Stable Inverted Perovskite
  Solar Cells. \emph{Angewandte Chemie International Edition} (2026).
  \url{https://doi.org/10.1002/anie.3095986}
\item
  Pallotta, R. et al.~Static antisolvent processing for scalable and
  sustainable perovskite photovoltaics. \emph{Communications Materials}
  (2026). \url{https://doi.org/10.1038/s43246-026-01321-8}
\item
  Zhang, T. et al.~Synergistic Crystallization Regulation by Substituent
  Electronic Effects of Polyfluoroarenes Toward Large Area Perovskite
  Solar Module Fabrication. \emph{Advanced Materials} (2026).
  \url{https://doi.org/10.1002/adma.74469}
\item
  Qi, J. et al.~Geometrically Asymmetric Phosphonium Zwitterions Enable
  Efficient Printable Mesoscopic Perovskite Solar Cells. \emph{Advanced
  Materials} (2026). \url{https://doi.org/10.1002/adma.74629}
\item
  Vaitukaitytė, D. et al.~Anchor Engineering of Hole‐Selective
  Enamine‐Based Self‐Assembling Monolayers for Perovskite Solar Cells.
  \emph{Advanced Science} (2026).
  \url{https://doi.org/10.1002/advs.76950}
\item
  Xu, Y. et al.~Closed‐Loop Slit‐Flushing Antisolvent Strategy Enables
  Controlled Three‐Dimensional Mass Transfer for Large‐Area Perovskite
  Solar Cells. \emph{Small} (2026).
  \url{https://doi.org/10.1002/smll.75068}
\item
  Kang, R. et al.~Thiourea Derivatives for Mitigating
  Photodegradation-Mediated Lead Leakage in Sustainable Perovskite Solar
  Cells and Modules. \emph{ACS Energy Letters} (2026).
  \url{https://doi.org/10.1021/acsenergylett.6c02033}
\item
  Wang, Z. et al.~Synergistic Self-Passivation of Vacancy Complexes in
  Metal Halide Perovskites. \emph{ACS Nano} (2026).
  \url{https://doi.org/10.1021/acsnano.6c10992}
\item
  Ouyang, D. et al.~Diminishing Defect and Ion Migration in
  Perovskite‐Organic Tandem Solar Cells by Ordered Anion‐Cation
  Modification Approach. \emph{Advanced Functional Materials} (2026).
  \url{https://doi.org/10.1002/adfm.77855}
\item
  Gao, Z. et al.~Selenocyanate Driven Organic Cations Homogenization in
  Wide‐Bandgap Perovskite. \emph{Advanced Energy Materials} (2026).
  \url{https://doi.org/10.1002/aenm.71462}
\item
  Chu, H. et al.~Polyfluorinated‐Regulator‐Assisted Scaffold‐Directed
  Tilted 1D/3D Heterointerfaces for High‐Voltage and Stable Perovskite
  Solar Cells. \emph{Angewandte Chemie International Edition} (2026).
  \url{https://doi.org/10.1002/anie.8256065}
\item
  Rout, P. et al.~Halide Exchange via Perovskite Nanocrystals for
  Interfacial Defect Passivation in Solar Cells. \emph{ChemRxiv} (2026).
  \url{https://doi.org/10.26434/chemrxiv.15007814/v1}
\item
  Huang, J. et al.~Sustainable Stabilization of Perovskite Solar Cells
  Using Multifunctional Ammonium Alginate for over 26\% Efficiency.
  \emph{Advanced Functional Materials} (2026).
  \url{https://doi.org/10.1002/adfm.78067}
\item
  Zhang, X. et al.~Multiple Passivation on High-Quality Perovskite
  Crystals Based on Low-Cost Polyfunctional Food Additive for
  High-Performance Perovskite Solar Cells in Ambient Air. \emph{ACS
  Applied Materials \& Interfaces} (2026).
  \url{https://doi.org/10.1021/acsami.6c10456}
\item
  Wang, C.-. et al.~Sulfamethoxazole-Induced Synergistic Interfacial
  Passivation for Efficient and Stable Perovskite Solar Cells. \emph{The
  Journal of Physical Chemistry Letters} (2026).
  \url{https://doi.org/10.1021/acs.jpclett.6c02253}
\item
  Duan, D. et al.~Ionic Liquid Engineering Induced Hydrogen‐Bond
  Networks for Efficient Wide‐Bandgap Perovskite Photovoltaics.
  \emph{Solar RRL} (2026). \url{https://doi.org/10.1002/solr.70461}
\item
  Chen, L. et al.~Dual‐Action Sulfonium‐Amino Acid Additive Enabling
  Crystallization Regulation and Defect Passivation for High‐Performance
  Carbon‐Electrode Printable Mesoscopic Perovskite Solar Cells.
  \emph{Angewandte Chemie International Edition} (2026).
  \url{https://doi.org/10.1002/anie.5667084}
\item
  Liu, X. et al.~Novel Multi-Pyridyl-Functionalized Spirobifluorenes as
  Precursor Additives for Improving Efficient and Stable Perovskite
  Solar Cells. \emph{ACS Applied Materials \& Interfaces} (2026).
  \url{https://doi.org/10.1021/acsami.6c13903}
\item
  Lin, S. et al.~Steric Hindrance‐Assisted Coordination for Efficient
  and Stable Tin‐Based Perovskite Solar Cells. \emph{Solar RRL} (2026).
  \url{https://doi.org/10.1002/solr.70454}
\item
  Liang, J. et al.~Solid-solution-induced lattice hardening reinforces
  perovskite stiffness for stable solar cell. \emph{National Science
  Review} (2026). \url{https://doi.org/10.1093/nsr/nwag484}
\item
  Zhao, J. et al.~Dynamic Strain Homogenization Regulation for
  Thermomechanically Robust Sn─Pb Perovskite Solar Cells. \emph{Advanced
  Energy Materials} (2026). \url{https://doi.org/10.1002/aenm.71497}
\item
  Sun, W. et al.~Competitive Ligand Modulation Reduces Buried
  Interfacial Voids and Defects Enabling 26.17\% Efficient Perovskite
  Solar Cells. \emph{Small} (2026).
  \url{https://doi.org/10.1002/smll.75360}
\item
  Zha, L. et al.~A Buried Sulfonate Molecular Bridge for Synchronous
  Charge Transport and Defect Passivation in High-Performance Perovskite
  Solar Cells. \emph{Nano-Micro Letters} (2026).
  \url{https://doi.org/10.1007/s40820-026-02278-6}
\item
  Guo, H. et al.~Molecular Disaggregation‐Bridging Engineered SnO 2
  Interfaces for Scalable Perovskite Photovoltaics. \emph{Advanced
  Energy Materials} (2026). \url{https://doi.org/10.1002/aenm.71419}
\item
  Hu, C. et al.~Hydrogen-bond network-reinforced self-assembled
  monolayer enables high efficiency and operational stability in
  pure-iodide wide-bandgap perovskite solar cells. \emph{DeCarbon}
  (2026). \url{https://doi.org/10.1016/j.decarb.2026.100190}
\item
  Zhang, W. et al.~Dipole-engineered co-adsorbed self-assembled
  monolayers boost Cs-based inorganic Pb-Sn perovskite tandem efficiency
  beyond 20\%. \emph{Nature Communications} (2026).
  \url{https://doi.org/10.1038/s41467-026-76847-z}
\item
  Zeng, X. et al.~Dipole-Moment-Tuned Self-Assembled Monolayers for
  Efficient Perovskite Solar Cells. \emph{ACS Applied Materials \&
  Interfaces} (2026). \url{https://doi.org/10.1021/acsami.6c11467}
\item
  Wang, H. et al.~Dual Modulation of the Buried Interface in Inverted
  Perovskite Solar Cells With Heteroatom‐Conjugated Asymmetric
  Self‐Assembled Molecules. \emph{Advanced Functional Materials} (2026).
  \url{https://doi.org/10.1002/adfm.77712}
\item
  Qin, X. et al.~Molecular Regulation toward SAM/Perovskite Interface
  for Efficient p‐i‐n Perovskite Solar Cells. \emph{Advanced Energy
  Materials} (2026). \url{https://doi.org/10.1002/aenm.71482}
\item
  Yang, J. et al.~Efficient and Stable Inverted Perovskite Solar Cells
  Via a Multi‐Arm Donor--Acceptor Dipole Molecular Bridge.
  \emph{Angewandte Chemie International Edition} (2026).
  \url{https://doi.org/10.1002/anie.3716107}
\item
  Yuan, S. et al.~Molecularly Templated Buried Interfaces for Inverted
  Perovskite Solar Cells. \emph{Angewandte Chemie International Edition}
  (2026). \url{https://doi.org/10.1002/anie.6932848}
\item
  Lu, C. et al.~Near-Surface Sulfidation-Induced p-n Junction Enabling
  High-Efficiency Inverted Perovskite Solar Cells. \emph{ACS Applied
  Materials \& Interfaces} (2026).
  \url{https://doi.org/10.1021/acsami.6c07884}
\item
  Zhu, Y. et al.~Bifunctional NiO x /Perovskite Interface Engineering
  for Efficient Air-Processed Wide-Bandgap Perovskite Solar Cells.
  \emph{ACS Applied Energy Materials} (2026).
  \url{https://doi.org/10.1021/acsaem.6c01708}
\item
  Li, Y. et al.~Structural Compensation Strategy: Composite NiO x Hole
  Transport Layers Enabling High‐Efficiency and Scalable Inverted
  Perovskite Solar Cells. \emph{Small} (2026).
  \url{https://doi.org/10.1002/smll.75056}
\item
  Ding, S. et al.~Inhibiting Buried Mechanical Failure via Vapor‐Induced
  Chemical Reconstruction for Durable Solar Cells. \emph{Small} (2026).
  \url{https://doi.org/10.1002/smll.75122}
\item
  Wang, X. et al.~Clean Tin Oxide via Volatile Additive Engineering for
  High‐performance Perovskite Photovoltaics. \emph{Advanced Energy
  Materials} (2026). \url{https://doi.org/10.1002/aenm.71514}
\item
  Zhou, S. et al.~The Impact of C60-Self-Assembled Monolayer
  Electron-Transport Layers in Negative--Intrinsic--Positive Perovskite
  Solar Cells. \emph{ACS Energy Letters} (2026).
  \url{https://doi.org/10.1021/acsenergylett.6c01485}
\item
  Bi, L. et al.~Crystallization Modulation for Stable Wide‐Bandgap
  Perovskite Solar Cells and Modules. \emph{Advanced Materials} (2026).
  \url{https://doi.org/10.1002/adma.74388}
\item
  Ali, W. et al.~Suppressing Non‐Radiative Recombination in Perovskite
  Solar Cells Through Integrated Tuning of Out‐of‐Plane Lattice Texture
  and Strain. \emph{Advanced Energy Materials} (2026).
  \url{https://doi.org/10.1002/aenm.71474}
\item
  Jiang, J. et al.~Macroporous PbI2 template and dual-mode coordination
  enable efficient perovskite solar cells. \emph{Journal of Colloid and
  Interface Science} (2026).
  \url{https://doi.org/10.1016/j.jcis.2026.141449}
\item
  Liu, T. et al.~Fluorinated Phenethylammonium and Pre-Annealing Static
  Stage Enable Intermediate-Mediated Stepwise Crystallization in Tin
  Perovskite Solar Cells. \emph{ACS Applied Materials \& Interfaces}
  (2026). \url{https://doi.org/10.1021/acsami.6c11057}
\item
  Kitamura, T. et al.~Solvent-Dependent Role of SnF2 in DMSO-Free Tin
  Halide Perovskite Solar Cells. \emph{ACS Applied Energy Materials}
  (2026). \url{https://doi.org/10.1021/acsaem.6c02063}
\item
  Yin, Z. et al.~Antiphase Boundaries Regulate Phase Stability and
  Performance in DMA + ‐Assisted CsPbI 3 ‐Based Perovskites.
  \emph{Advanced Materials} (2026).
  \url{https://doi.org/10.1002/adma.74538}
\item
  Mikuciunaite, J. et al.~Fluorinated Methylammonium Cation Containing
  Perovskite Solar Cells With Over 25\% Power Conversion Efficiency.
  \emph{Advanced Energy Materials} (2026).
  \url{https://doi.org/10.1002/aenm.71396}
\item
  Yang, Y. et al.~Precise Management of n‐Value Distribution at the
  2D/3D Interface Enabling Efficient Wide‐Bandgap Perovskite Solar
  Cells. \emph{Advanced Science} (2026).
  \url{https://doi.org/10.1002/advs.76791}
\item
  Thor, W. et al.~Diketopyrrolopyrrole‐Based Cation as a Semiconducting
  Spacer in Layered Perovskite. \emph{Advanced Materials} (2026).
  \url{https://doi.org/10.1002/adma.74681}
\item
  Jouane, Y. et al.~Stable lead-free bismuth perovskite solar cells via
  thermally activated bilayer conversion. \emph{Solar Energy} (2026).
  \url{https://doi.org/10.1016/j.solener.2026.114961}
\item
  Glazer, K. et al.~Site-selective alkali ion alloying in double
  perovskites via liquid-assisted mechanosynthesis revealed by
  solid-state NMR. \emph{ChemRxiv} (2026).
  \url{https://doi.org/10.26434/chemrxiv.15008057/v1}
\item
  Padova, P.D. et al.~Eco-friendly water-synthesis of lead-free
  Cs2AgBiBr6 perovskites. \emph{Communications Materials} (2026).
  \url{https://doi.org/10.1038/s43246-026-01261-3}
\item
  Gong, H. et al.~Rationally Designed Cooperative Self‐Assembled
  Monolayers Enable Holistic Modulations of Wide‐Bandgap Perovskites for
  High‐Performance Perovskite‐Silicon Tandem Solar Cells. \emph{Advanced
  Energy Materials} (2026). \url{https://doi.org/10.1002/aenm.71512}
\item
  Park, S.J. et al.~Self-Assembled Monolayers for High-Efficiency
  Perovskite/Si Tandem Solar Cells. \emph{ACS Energy Letters} (2026).
  \url{https://doi.org/10.1021/acsenergylett.6c01966}
\item
  Mahmoud, M.A.A. et al.~Impact of silicon substrate topography on
  sequentially evaporated perovskite film growth. \emph{Joule} (2026).
  \url{https://doi.org/10.1016/j.joule.2026.102646}
\item
  Zhao, Z. et al.~High‐Performance Perovskite/Silicon Tandem Solar Cells
  via Ion Modulation. \emph{Advanced Functional Materials} (2026).
  \url{https://doi.org/10.1002/adfm.77637}
\item
  Zhang, T. et al.~Dimensional Seed Engineering of Strain in
  Wide-Bandgap Perovskites for Efficient and Stable Tandem Solar Cells.
  \emph{ACS Nano} (2026). \url{https://doi.org/10.1021/acsnano.6c07355}
\item
  Yuan, J. et al.~Universal Sacrificial Coordination Strategy for
  ALD‐Resilient SAMs Achieving High‐Performance Perovskite/Organic
  Tandem Solar Cells. \emph{Advanced Materials} (2026).
  \url{https://doi.org/10.1002/adma.74461}
\item
  Bristow, H. et al.~Material Selection through Accelerated Aging
  Enables \textgreater95\% Performance Retention of a
  Perovskite--Silicon Tandem Solar Cell after 1 Year of Outdoor
  Monitoring. \emph{ACS Energy Letters} (2026).
  \url{https://doi.org/10.1021/acsenergylett.6c01363}
\item
  Li, L. et al.~Large‐Scale Perovskite Photovoltaics: Mitigating Defect
  Challenges in Upscaled Modules and Monolithic Tandem Cells.
  \emph{Advanced Functional Materials} (2026).
  \url{https://doi.org/10.1002/adfm.202532038}
\item
  Zhang, J. et al.~Decoupling Bidirectional Photochemical Degradation
  via Radical Scavenging for Stable Inverted Perovskite Solar Cells.
  \emph{Angewandte Chemie International Edition} (2026).
  \url{https://doi.org/10.1002/anie.4221605}
\item
  Yan, X. et al.~Self-assembled monolayers with cross-linkable header on
  the linker for durable perovskite photovoltaics. \emph{Science
  Advances} (2026). \url{https://doi.org/10.1126/sciadv.aec3096}
\item
  Luo, C. et al.~Symmetry-breaking co-assembly of conjugated molecules
  boosts perovskite photovoltaics. \emph{Nature Communications} (2026).
  \url{https://doi.org/10.1038/s41467-026-76848-y}
\item
  Guo, Y. et al.~Regulating Molecular Adsorption Configuration via
  Spatially Balanced Polarity for High‐Performance Inverted Perovskite
  Solar Cells. \emph{Advanced Functional Materials} (2026).
  \url{https://doi.org/10.1002/adfm.77942}
\item
  Wang, K. et al.~Acid‐Base Complexation Induced Dipole Engineering for
  Durable Inverted Perovskite Photovoltaics. \emph{Advanced Materials}
  (2026). \url{https://doi.org/10.1002/adma.74438}
\item
  Tang, Z. et al.~Mitigating Stress‐Induced Nonphotoactive Phase
  Transition Through Sodium Sulfonate Engineering for Stable and
  Efficient Perovskite Solar Cells. \emph{Angewandte Chemie
  International Edition} (2026).
  \url{https://doi.org/10.1002/anie.2425823}
\item
  Ma, X. et al.~Regulating lattice evolution via an ordered
  two-dimensional intermediate phase for resilient perovskite solar
  cells. \emph{Nature Communications} (2026).
  \url{https://doi.org/10.1038/s41467-026-76425-3}
\item
  Abbas, Z. et al.~Suppression of Lattice Heterogeneity for Efficient
  Inverted Perovskite Solar Cells. \emph{Advanced Functional Materials}
  (2026). \url{https://doi.org/10.1002/adfm.77729}
\item
  Lukas, T. et al.~Beyond the Gold Standard: Towards Industrially Viable
  Electrodes for Durable Perovskite Solar Cells. \emph{Advanced Energy
  Materials} (2026). \url{https://doi.org/10.1002/aenm.71372}
\item
  Liu, N. et al.~Two Birds With One Stone: Modifier Interfacial
  Engineering for High‐Efficiency Stable p‐i‐n All‐Inorganic
  Perovskite/Organic Tandem Solar Cells. \emph{Advanced Functional
  Materials} (2026). \url{https://doi.org/10.1002/adfm.77805}
\item
  Ren, H. et al.~Pushing the Interconnection Limit in Perovskite‐Organic
  Tandem Solar Cells Through Recombination Channels Delocalization.
  \emph{Angewandte Chemie International Edition} (2026).
  \url{https://doi.org/10.1002/anie.6548693}
\end{enumerate}

\end{document}
